\chapter{到達目標}

\begin{figure}[htbp]
 \centering
 \includegraphics[width=3cm]{images/apple.eps}
 \caption{リンゴ}
 \label{fig:apple}
\end{figure}

\section{本プロジェクトにおける目的}\label{sec:mokuteki}
%1.3節で述べた課題をより具体的に記述する。成果に対して必ず満たすべき条件を含む

地域の特色を生かしたおいしいカレーの作り方が課題。
最終的には、100人中75人以上がおいしいというカレーの詳細なレシピを作ること。
またそのカレーは函館の海産物を用いたものであること。
レシピの手順は、なぜその方法がいいのかも含めて記述されていること。
\bunseki{未来}

\subsection{通常の授業ではなく、プロジェクト学習で行う利点}

本課題では材料に多種多様なものが考えられるが、複数の人数で試作することにより、 
さまざまなバリエーションが試せる。また、味の好みの偏りが少なくなる。 
通常の授業では基本的に個人の知識・技術について講義・演習形式で行われるため、 
共同作業で行うべき作業時間の多いテーマに関しては向かない。

\bunseki{函館}

\subsection{地域との関連性（必要ならば）}

海産物を特徴にしたカレーができると地域の名物料理として売り出せるかも。 
また函館の特産品の売り上げが伸びるかも。

\bunseki{北海}


\section{具体的な手順・課題設定}\label{sec:tejun}
%2.1節で述べた課題を解決するための小課題を手順に分け、その詳細を記述する。 各人への割り当て可能なレベルまで具体化する。なお、以下を必ず含むこと。 
%・このような課題設定に至るプロセス
%・各小課題の解決過程に関連する講義 
%・各小課題の解決過程で用いる既存技術、また習得技術

数多いレシピをもとに効率的に函館特産物を用いたカレーを製作し、コスト面から、 
試作の数を20種類以内に収める目的で、情報収集に力をいれ、以下のように 
手順を設定した。

\begin{enumerate}
\item 従来のカレーレシピ収集（料理本・テレビ・Web)
\par 課題：レシピを共通する部分と異なる部分にわけ、グループ化する。
異なる部分については、それぞれのメリットデメリットを挙げる。

\item 函館特産食品の種類の調査
\par 課題：生産高が多く、一般に特産品として知名度の高いものを調査する。
季節・標準的な値段・一般的な調理法とその調理法により引き出せる味の調査をする。

\item 函館特産食品とカレーとの組み合わせを調べる（過去のレシピの検索）
\par 課題：省略

\item 各材料の下ごしらえのレシピ化。
\par 課題：省略

\item 試作レシピパターンの決定
\par 課題：省略

\item 試作
\par 課題：省略

\item アンケート実施及び解析、改善点の発見
\par 課題：省略

\item 好評な試作パターンについてのバリエーションを設定。
\par 課題：省略

\item アンケート実施及び解析、改善点の発見
（75パーセント以上の好評価を得るまで8-9の繰り返し）
\par 課題：省略
\end{enumerate}

\bunseki{未来}

\section{課題の割り当て}
%2.2節で具体化した各小課題を誰にどのように分担したか、またその理由も含めて記述する。

各人の得意分野及び関連性、時間軸のスケジュールを基準に
以下のように割り当てた。
\bunseki{函館}
